% GENERATED FILE. Edit canonical lessons, then run npm run generate:books.
% canonical-source-hash: fnv1a64:28dd1822c966c53c
% canonical-lessons: LA-C448-sterto, LA-C448-strepo, LA-C448-obstrepo, LA-C448-potor, LA-C448-raucus

\chapter{To snore, To make a noise, To make a din, Drinker, Hoarse}
\label{ch:la-to-snore-to-make-a-noise-to-make-a-din-drinker-hoarse}
\hlchaptermodality{448}

\begin{chapteropening}
\textbf{By the end of this chapter:} \emph{I can complain about a noisy night: someone snoring, people making noise, a din, the drinkers and their hoarse voices.}

Complete the last lesson of chapter 448: I can complain about a noisy night: someone snoring, people making noise, a din, the drinkers and their hoarse voices.
\end{chapteropening}

\section[stertō, stertere]{stertō, stertere — to snore}

\label{lesson:LA-C448-sterto}



\begin{quote}
Before the new one: say the Latin for sleepless, then the Latin for to lie down.
\end{quote}

\subsection*{You'll want to know: stertō, stertere}
\textbf{stertō, stertere} — \textquotedblleft{}to snore\textquotedblright{}.

\begin{quote}
Hospes ēbrius tōtam noctem stertit.
\end{quote}

\begin{quote}
\emph{The drunken guest snores all night.}
\end{quote}

\emph{stertō, stertere, stertuī}. English \emph{stertorous} 'heavy and snoring' comes from it.

\begin{cousinweb}
The origin of \emph{stertō} is not settled: Lewis \& Short compare Greek \emph{darthanō} 'to sleep', but EtymDB gives no parent. English \emph{stertor} and \emph{stertorous} come from it.
\end{cousinweb}

\begin{grammarlens}[title={What you've built}]
One more word for telling, explaining and describing at B1.
\end{grammarlens}

\subsection*{Your turn}
\begin{itemize}
\raggedright
  \item \emph{Say it:} \emph{stertō, stertere}
  \item \emph{Say it:} \emph{stertō, stertere}, once more
  \item \emph{Say it:} a sentence of your own with \emph{stertō, stertere}
  \item \emph{Recall:} say the Latin for sleepless, then the Latin for to lie down, then say \emph{stertō, stertere} again
\end{itemize}

\begin{tcolorbox}[breakable,colback=teal!4,colframe=teal!35!black,title={Before you move on}]
What does stertō, stertere mean? (\textquotedblleft{}To snore\textquotedblright{}.) Say it once more.
\end{tcolorbox}
\section[strepō, strepere]{strepō, strepere — to make a noise, to clatter, to roar}

\label{lesson:LA-C448-strepo}



\begin{quote}
Before the new one: say the Latin for to lie down, then the Latin for to snore.
\end{quote}

\subsection*{You'll want to know: strepō, strepere}
\textbf{strepō, strepere} — \textquotedblleft{}to make a noise, to clatter, to roar\textquotedblright{}.

\begin{quote}
In viā plaustra strepunt, et nēmō dormīre potest.
\end{quote}

\begin{quote}
\emph{In the street the wagons clatter, and nobody can sleep.}
\end{quote}

\emph{strepō, strepere, strepuī, strepitum}. Any loud noise: wheels, voices, water.

\begin{cousinweb}
The origin of \emph{strepō} is not settled: Lewis \& Short give its senses, to rattle, rumble, roar, but no etymology, and EtymDB gives no parent.
\end{cousinweb}

\begin{grammarlens}[title={What you've built}]
One more word for telling, explaining and describing at B1.
\end{grammarlens}

\subsection*{Your turn}
\begin{itemize}
\raggedright
  \item \emph{Say it:} \emph{strepō, strepere}
  \item \emph{Say it:} \emph{strepō, strepere}, once more
  \item \emph{Say it:} a sentence of your own with \emph{strepō, strepere}
  \item \emph{Recall:} say the Latin for to lie down, then the Latin for to snore, then say \emph{strepō, strepere} again
\end{itemize}

\begin{tcolorbox}[breakable,colback=teal!4,colframe=teal!35!black,title={Before you move on}]
What does strepō, strepere mean? (\textquotedblleft{}To make a noise, to clatter, to roar\textquotedblright{}.) Say it once more.
\end{tcolorbox}
\section[obstrepō, obstrepere]{obstrepō, obstrepere — to make a din, to drown out with noise}

\label{lesson:LA-C448-obstrepo}



\begin{quote}
Before the new one: say the Latin for to snore, then the Latin for to make a noise.
\end{quote}

\subsection*{You'll want to know: obstrepō, obstrepere}
\textbf{obstrepō, obstrepere} — \textquotedblleft{}to make a din, to drown out with noise\textquotedblright{}.

\begin{quote}
Virī ēbriī in caupōnā obstrepunt, et hospitēs queruntur.
\end{quote}

\begin{quote}
\emph{Drunken men make a din in the inn, and the guests complain.}
\end{quote}

\emph{obstrepō, obstrepere, obstrepuī, obstrepitum}, often with a dative: to roar at someone. English \emph{obstreperous} comes from it.

\begin{cousinweb}
\emph{obstrepō} is built from \emph{ob-} 'against' and \emph{strepō} 'to make a noise' (Lewis \& Short). English \emph{obstreperous} comes from \emph{obstreperus}, built on it.
\end{cousinweb}

\begin{grammarlens}[title={What you've built}]
One more word for telling, explaining and describing at B1.
\end{grammarlens}

\subsection*{Your turn}
\begin{itemize}
\raggedright
  \item \emph{Say it:} \emph{obstrepō, obstrepere}
  \item \emph{Say it:} \emph{obstrepō, obstrepere}, once more
  \item \emph{Say it:} a sentence of your own with \emph{obstrepō, obstrepere}
  \item \emph{Recall:} say the Latin for to snore, then the Latin for to make a noise, then say \emph{obstrepō, obstrepere} again
\end{itemize}

\begin{tcolorbox}[breakable,colback=teal!4,colframe=teal!35!black,title={Before you move on}]
What does obstrepō, obstrepere mean? (\textquotedblleft{}To make a din, to drown out with noise\textquotedblright{}.) Say it once more.
\end{tcolorbox}
\section[pōtor, pōtōris]{pōtor, pōtōris — a drinker, a heavy drinker}

\label{lesson:LA-C448-potor}



\begin{quote}
Before the new one: say the Latin for to make a noise, then the Latin for to make a din.
\end{quote}

\subsection*{You'll want to know: pōtor, pōtōris}
\textbf{pōtor, pōtōris} — \textquotedblleft{}a drinker, a heavy drinker\textquotedblright{}.

\begin{quote}
In tabernā pōtōrēs multum vīnum bibunt.
\end{quote}

\begin{quote}
\emph{In the tavern the drinkers drink a lot of wine.}
\end{quote}

Third declension, masculine: \emph{pōtor, pōtōris m.}, from the same root as \emph{pōtus} 'drunk'.

\begin{cousinweb}
\emph{pōtor} is built on the root of \emph{pōtō} 'to drink' with the suffix \emph{-tor} 'one who' (Lewis \& Short). Greek \emph{pinō} 'to drink' is a cousin of that root.
\end{cousinweb}

\begin{grammarlens}[title={What you've built}]
One more word for telling, explaining and describing at B1.
\end{grammarlens}

\subsection*{Your turn}
\begin{itemize}
\raggedright
  \item \emph{Say it:} \emph{pōtor, pōtōris}
  \item \emph{Say it:} \emph{pōtor, pōtōris}, once more
  \item \emph{Say it:} a sentence of your own with \emph{pōtor, pōtōris}
  \item \emph{Recall:} say the Latin for to make a noise, then the Latin for to make a din, then say \emph{pōtor, pōtōris} again
\end{itemize}

\begin{tcolorbox}[breakable,colback=teal!4,colframe=teal!35!black,title={Before you move on}]
What does pōtor, pōtōris mean? (\textquotedblleft{}A drinker, a heavy drinker\textquotedblright{}.) Say it once more.
\end{tcolorbox}
\section[raucus, rauca, raucum]{raucus, rauca, raucum — hoarse; harsh-sounding}

\label{lesson:LA-C448-raucus}



\begin{quote}
Before the new one: say the Latin for to make a din, then the Latin for a drinker.
\end{quote}

\subsection*{You'll want to know: raucus, rauca, raucum}
\textbf{raucus, rauca, raucum} — \textquotedblleft{}hoarse; harsh-sounding\textquotedblright{}.

\begin{quote}
Post noctem longam vōx caupōnis rauca est.
\end{quote}

\begin{quote}
\emph{After a long night the innkeeper's voice is hoarse.}
\end{quote}

An adjective of three endings. English \emph{raucous} comes from it.

\begin{cousinweb}
\emph{raucus} is built on \emph{rāvis} 'hoarseness' with an adjective suffix (EtymDB); Lewis \& Short trace it to a root meaning 'to make a loud noise'. English \emph{raucous} comes from it.
\end{cousinweb}

\begin{grammarlens}[title={What you've built}]
One more word for telling, explaining and describing at B1.
\end{grammarlens}

\subsection*{Your turn}
\begin{itemize}
\raggedright
  \item \emph{Say it:} \emph{raucus, rauca, raucum}
  \item \emph{Say it:} \emph{raucus, rauca, raucum}, once more
  \item \emph{Say it:} a sentence of your own with \emph{raucus, rauca, raucum}
  \item \emph{Recall:} say the Latin for to make a din, then the Latin for a drinker, then say \emph{raucus, rauca, raucum} again
\end{itemize}

\begin{tcolorbox}[breakable,colback=teal!4,colframe=teal!35!black,title={Before you move on}]
What does raucus, rauca, raucum mean? (\textquotedblleft{}Hoarse; harsh-sounding\textquotedblright{}.) Say it once more.
\end{tcolorbox}
